% !TeX root = ../papers/nc-p3.tex
\section{過剰曲線の三次評価による全八余因子窓の排除}
\label{sec:eight-cofactor-excess}
\begin{theorem}\label{thm:eight-cofactor-excess}
四辺有限の標準有限Frobenius窓では、整な反対八次曲面の余因子の包が八次元となることはない。
非被約な点窓と特異核の全鉛筆元を含む。従って整な場合の余因子数は $7,9,10$ に限られる。
\end{theorem}
\begin{proof}[水平成分がない場合]
滑らかな核と特異境界は既に排除した。前節の特異有限鉛筆の $M,Y,F,H,j$ を用いる。
$Y$ は入力の全階数1軌跡なので $X$ を含み、次元は高々1である。
$j=0$ ならスカラー束を完備化した $E$ 上の線織曲面で四行三列の退化は有限である。
垂直曲線なら無限遠と交わり、そこでは包含が単射である。非分岐では次数零の線束と自明束の直和、
分岐では自明束同士の拡張を射影化すればよい。源は $A\otimes\mathcal O(-1)$、
的の四成分には底の次数零の捩りだけが現れる。二次の対称代数のKoszul複体は終端余核を除いて完全である。
実際有限集合外では全射な双対写像の核が一本で、次数2では完全、また長さ2の自由複体であり
底曲面はCohen--Macaulayである。終端余核の長さは16である。
$\deg A^*=4,\ \deg\operatorname{Sym}^2A^*=16$ と、相対 $\mathcal O(1),\mathcal O(2)$ の直像が
階数2,3・次数零であることから、種数一のEuler数で $3\cdot16-4\cdot2\cdot4=16$ を得る。
定スカラー・定残余点のファイバーは入力点または直線である。平面なら整な $D$ に平面が入る。
$X$ の直線切断長は3以下、点の切断長は1以下なので相対二次式が各ファイバーの切断を生成する。
中山の補題で $\ell(X)\le16$、矛盾する。
\end{proof}

\begin{proof}[曲線成分を含む零スキームのEuler数]
$j>0$ とする。$\mathcal O_Y$ の最大有限長部分を除いたCohen--Macaulay曲線を $C$、
$b=\deg_C H$ とする。$C$ は一般重複度込みで底を次数 $j$ で覆う水平部分と垂直部分からなる。
三束 $F$ の切断のKoszul複体は $\mathcal O_Y$ と第一ホモロジーだけを持ち、後者は
\[
               \omega_C\otimes(\omega_M\otimes\det F)^{-1}
\]
となる。実際零軌跡の余次元は2で、完全性の深さ判定で第二以上のホモロジーは零。
Koszul複体を $\omega_M$ で双対化し次数2を読むと、余次元2に支えを持つ第一ホモロジーの
零次・一次Extは寄与せず、$\mathcal E xt^2(\mathcal O_Y,\omega_M)=\omega_C$ を得る。
三次元Gorensteinの $M$ では有限長部分の第二Extも零である。
$\omega_M=-3H+\pi^*\det A^*$、$\deg\det A^*=4$、$\det F$ は数値的に $4H$ である。
Koszul複体の捩りに依らない交代Euler数は $c_3(F)=4H^3=16$ である。
これは切断の正則性を必要とせず、束のKoszul項の交代Chern指標の最初の非零次数が3となることによる。
従って曲線の双対性とEuler数の次数公式から
\[
       \chi(\mathcal O_Y(3H))=16+2b-4j-\chi(\mathcal O_C).                 \tag{1}
\]

一般の滑らかな $z\in E$ で、$a(z)=b'(z)=0$ の入力直線を $\ell_z$ とする。
$b'$ は第二切断または一次持上げの残る局所成分を指す。この二条件が従属なら $D$ に平面が入る。
直線族は $D$ を覆う。実際各直線で $N$ は $z$ の評価を消し、四次式 $a$ が一般の全 $z$ で
消えることはないからこれらの直線は動く。従って一般の直線で制限定数核を除いた対称三行の
六次行列式は非零である。全二次成分に基点がないのでその直線上の各点の第一単因子は単元、
六次の零位数は階数1切断の局所長の二倍以上で、$\ell(Y\cap\ell_z)\le3$ である。

$C$ 上の $a(z)$ の切断は、$u=z$ の因子と共通三点因子を $z$ で評価した因子の積である。
滑らかな $z$ の近くでは点イデアル $I_u$ が対角のCartier式を持つため、この因数分解は冪零を含む。
一般の $z$ を選べば各因子は $C$ のどの成分にも恒等零でなく、零因子の長さは加法的である。
$u=z$ の部分は $j$、残りは $b-j$ であり、後者では $a,b'$ がともに $z$ で消える。
分岐元でも一定の三点因子が一次持上げで消えることから同じ結論となる。
$Y$ の入力への射が閉埋込みであるため、この残余長は入力直線で失われない。ゆえに
\[
                              b-j\le3 .                                     \tag{2}
\]
\end{proof}

\begin{proof}[垂直な非被約部分の切断]
水平部分のスキーム閉包を $C_h$ とする。前節の導手の固有値計算により、
$j=1$ なら $C_h\simeq\mathbb P^1$、$j=2$ なら二本の $\mathbb P^1$ が二節点で交わる被約曲線である。
前者のEuler数は1、後者は0である。入力 $H$ は各成分で正次数なのでいずれも $H^1(C_h,3H)=0$
（後者は二つの貼合点への評価を使う）。導手で水平枝と垂直直線は交わらない。
節点では二行の固有値が相異なるためである。尖点でも垂直直線が交われば、導手の二行の
自己準同型がその固有値のスカラー行列になる。しかし正規化局所変数 $r$ でその行列の一次項は零、
水平固有値の一次項は非零であり、固有ベクトルの方程式の一次項が非零となって矛盾する。

$T=\ker(\mathcal O_C\to\mathcal O_{C_h})$ は垂直直線上に支えられる。
各直線は残余点とスカラーの固定したファイバーである。これら二つをパラメータとする局所底の
極大イデアルを $\mathfrak n$ とする（導手では底の局所環自体を使う）。
このファイバーは線型な固有空間なので、$T$ の $\mathfrak n$ 冪による有限フィルトレーションの各商は
入力直線上の層となる。各商は $\mathcal O$ と $\mathcal O(-1)$ の有限直和の商である。
実際水平曲線が直線と交わらなければ構造層の生成元1で足りる。
交わる場合は導手外で、異なる水平スカラーは互いに異なるから近傍には水平切断が一本だけである。
底を自明化すればそのイデアルは底のスカラーの式と、固有直線を切る相対一次式で生成される。
$T$ はそのイデアルの $C$ 上の像であり、$\mathfrak n$ の各次数を掛けても相対生成次数は増えない。
従って各商は捩り0でも $3H$ でも第一コホモロジーが零であり、特に各Euler数は非負である。
この議論は垂直直線の重複度、水平切断からの接触次数を仮定しない。これで
\[
         \chi(\mathcal O_C)\ge2-j,\qquad H^1(C,3H)=H^1(Y,3H)=0
\]
を得る。点窓上で三次評価が単射なので $20\le h^0(Y,3H)$ が必要であるが、
(1),(2) の上界は $j=1$ なら19、$j=2$ なら18である。これで残る有限元を全て排除した。
\end{proof}

\path{computations/verify_eight_cofactor_excess.py} は二種類の長さ16の計算と水平重複度別の三次上界を確認する。
